%% T5.tex -- Proposition 5: NHPP limit under rare-failure scaling.
%% REVISED: separates survival S(d) from cumulative first-passage R(d);
%% proves the Goel--Okumoto limit for the CDF of T_oplus rather than for
%% the Bernoulli survival function, repairing the sign/monotonicity
%% inconsistency in the earlier draft.

\noindent\emph{The following is a direct instance of the classical
Le~Cam / Poisson limit theorem applied to a triangular array of
Bernoulli step-outcomes, adapted to yield the Goel--Okumoto NHPP
mean-value function~\cite{goelokt} as the limit of the
\emph{cumulative first-passage CDF} $R(d)$ (Def.~\ref{def:reliab}),
not the survival function. The mathematics is textbook; the
contribution is the \emph{micro-foundation} reading: modern LLM agents
re-emerge as classical Goel--Okumoto software reliability limits under
rare-failure scaling, licensing 40 years of NHPP estimation tooling for
the agent setting.}

\paragraph{Notation (survival vs.\ cumulative success).}
To eliminate an ambiguity we became aware of during revision, we
reserve the symbol $R(d)$ strictly for the \emph{cumulative
first-passage CDF} to the success state, i.e., the
monotone-\emph{increasing} function
\begin{equation}
R(d) \;\triangleq\; \Pr\!\bigl[T_\oplus \le d\bigr], \qquad R(0)=0,\;
R(d)\uparrow R_\infty \in [0,1],
\label{eq:R-def}
\end{equation}
where $T_\oplus = \inf\{t\ge 0: X_t=\oplus\}$ is the hitting time to
$\oplus$. We separately denote the \emph{survival function} (the
probability of \emph{non}-absorption to $\oplus$ by step $d$) by
\begin{equation}
S(d) \;\triangleq\; \Pr\!\bigl[T_\oplus > d\bigr],
\qquad S(0)=1,\; S(d)\downarrow 1-R_\infty.
\label{eq:S-def}
\end{equation}
In general $R(d)+S(d)=1$ along any trajectory that eventually absorbs
in $\oplus$; and $R(d)+S(d)+\Pr[T_\ominus\le d]$ accounts for the full
mass when $\ominus$ is also reachable.

\begin{proposition}[NHPP Rare-Failure Limit; adapted from Goel--Okumoto
\cite{goelokt}]
\label{thm:nhpp}
Let $(\mathcal{M}_n)_{n\ge 1}$ be a sequence of agent Markov chains,
each with a single transient state $s_0$, per-step success probability
$\mu_n$, and per-step failure probability $\varepsilon_n$, so that
$Q_n = 1-\mu_n-\varepsilon_n$,
$R_{\oplus,n}=\mu_n$, $R_{\ominus,n}=\varepsilon_n$.  Assume the
\emph{rare-failure} scaling
\begin{equation}
\mu_n\to 0,\quad \varepsilon_n\to 0,\quad
d_n\mu_n\to \Lambda\in(0,\infty),\quad
\varepsilon_n/\mu_n\to 0.
\label{eq:t5-scaling}
\end{equation}
Then the cumulative first-passage CDF $R_n(d)$ (from \eqref{eq:R-def})
converges, for every fixed $c\in(0,\infty)$ and $d=cd_n$, to
\begin{equation}
R_n(d)\;\xrightarrow[n\to\infty]{}\; 1-e^{-c\Lambda},
\label{eq:t5-GO}
\end{equation}
which is the Goel--Okumoto NHPP mean-value function
$m(t)=a(1-e^{-bt})$ with $a=1$, $b=\Lambda/d_n$. In particular,
$R_n$ is \emph{monotone increasing} in $d$, $R_n(0)=0$, and
$R_n(d)\to 1$ as $d\to\infty$.  The step-level failure count
$F_n$ also converges: $F_n\xrightarrow{d}\mathrm{Poisson}(0)$ (zero in
the limit), which is the usual rare-failure regime.
\end{proposition}

%% ----------------------------------------------------------------
%% [REDUCTION 2026-04-24] Proof and remarks suppressed from body.
%% ----------------------------------------------------------------
\iffalse
\begin{proof}
\emph{Step 1 --- exact CDF by the fundamental matrix.}
By Proposition~\ref{thm:closed} (T1) applied to the 1-state chain,
\begin{equation}
R_n(d) \;=\; \mu_n \sum_{t=0}^{d-1} Q_n^{\,t}
\;=\; \mu_n\,\frac{1-Q_n^{\,d}}{1-Q_n}
\;=\; \frac{\mu_n}{\mu_n+\varepsilon_n}\bigl(1-(1-\mu_n-\varepsilon_n)^d\bigr).
\label{eq:t5-exact-R}
\end{equation}
This is a CDF: it is $0$ at $d=0$ and approaches
$\mu_n/(\mu_n+\varepsilon_n)\uparrow 1$ as $d\to\infty$.  (The
corresponding survival function is
$S_n(d) = 1 - R_n(d) - \frac{\varepsilon_n}{\mu_n+\varepsilon_n}
         \bigl(1-(1-\mu_n-\varepsilon_n)^d\bigr)
       = (1-\mu_n-\varepsilon_n)^d$,
which is monotone \emph{decreasing} from $1$ to $0$. This is the
quantity that was mislabeled as $R(d)$ in earlier drafts; it is
the survival function of $T_\oplus$, not its CDF.)

\emph{Step 2 --- algebraic limits.}
Two factors in \eqref{eq:t5-exact-R}.
(a) The prefactor: since $\varepsilon_n/\mu_n\to 0$,
$\mu_n/(\mu_n+\varepsilon_n)\to 1$.
(b) The exponent:
$\log(1-\mu_n-\varepsilon_n) = -(\mu_n+\varepsilon_n)(1+o(1))$ as
$n\to\infty$, so
$(1-\mu_n-\varepsilon_n)^d = \exp(d\log(1-\mu_n-\varepsilon_n))
 = \exp(-d(\mu_n+\varepsilon_n)(1+o(1)))$. Using $d=cd_n$
and $d_n\mu_n\to\Lambda$, $d_n\varepsilon_n/\mu_n\to 0
\Rightarrow d_n\varepsilon_n\to 0$:
$d(\mu_n+\varepsilon_n) = cd_n(\mu_n+\varepsilon_n) \to c\Lambda.$
Combining (a) and (b), $R_n(d) \to 1-e^{-c\Lambda}$, establishing
\eqref{eq:t5-GO}.

\emph{Step 3 --- generalization to $m$-state forward chains.}
For a general forward chain with $m$ transient states and per-state
self-loop probability $1-\mu_n^{(i)}-\varepsilon_n^{(i)}$, the
hitting time $T_\oplus$ is a sum of $m$ independent geometric
random variables with rates $\mu_n^{(i)}+\varepsilon_n^{(i)}$.
Under \eqref{eq:t5-scaling} applied uniformly across states, a
generalized central-limit / Poisson-limit argument on the triangular
array of per-step Bernoulli indicators gives the same limit shape
$R(d)\to 1-e^{-\Lambda d/d_n}$ after a rescaling of time by $m$.
The resulting mean-value function is precisely Goel--Okumoto with
asymptotic reliability $a=1$ and rate $b=\Lambda/(md_n)$.

\emph{Step 4 --- step-level failure count.}
Under \eqref{eq:t5-scaling}, $d_n\varepsilon_n\to 0$, so by the
classical Poisson limit theorem for triangular arrays
(Le~Cam's theorem, $d_{\rm TV}(\mathrm{Bin}(d_n,\varepsilon_n),
\mathrm{Poisson}(d_n\varepsilon_n))\le \varepsilon_n\to 0$), the
number of step-level failures $F_n\xrightarrow{d}\mathrm{Poisson}(0)$,
i.e., failures vanish in the scaling limit, consistent with the
classical ``rare-failure'' regime of Goel--Okumoto.
\end{proof}

\begin{remark}[Scope and limitations]
Proposition~\ref{thm:nhpp} holds under rare-failure scaling
\eqref{eq:t5-scaling}.  For moderate failure rates ($\varepsilon
\gtrsim 0.1$) or comparable success/failure rates
($\varepsilon_n/\mu_n\not\to 0$), the limit \eqref{eq:t5-GO} does not
hold and the exact closed form from Proposition~\ref{thm:closed} should
be used instead.
\end{remark}

\begin{remark}[Micro-foundation for NHPP]
Classical NHPP software reliability models (Goel--Okumoto,
Musa--Iannino--Okumoto) postulate an intensity function $\lambda(t)$
for fault occurrences and fit it phenomenologically.
Proposition~\ref{thm:nhpp} shows that the Goel--Okumoto mean-value
function $1-e^{-bd}$ \emph{re-emerges from first principles} as the
limit of the cumulative first-passage CDF of a rare-failure agent
Markov chain.  The agent-chain formalism thus provides the previously
missing mechanistic justification for applying NHPP estimators to LLM
agent reliability.
\end{remark}

\begin{remark}[Re: Kolmogorov--Smirnov test in SS4]
Simulation study SS4 fits the Goel--Okumoto model to the empirical
first-passage CDF $R_n(d)=\Pr[T_\oplus\le d]$ (not the survival
function $S_n(d)$) for several values of $\varepsilon_n/\mu_n$ and
reports the two-sample KS $p$-value; the null hypothesis
(empirical CDF follows $1-e^{-\mu d}$) is retained at $\alpha=0.05$
for $\varepsilon_n/\mu_n\le 0.1$ and $\mu_n d_n\to\Lambda=2$ fixed.
\end{remark}
\fi
